\documentclass{article}
\input{../../../lib.sty}

\title{\texttt{fn make\_private\_group\_by}}
\author{Michael Shoemate}
\date{}

\begin{document}

\maketitle

\contrib
Proves soundness of \texttt{fn make\_private\_group\_by} in \asOfCommit{mod.rs}{0db9c6036}.

\section{Hoare Triple}
\subsection*{Precondition}
\subsubsection*{Compiler-verified}
\begin{itemize}
    \item Generic \texttt{MS} must implement trait \texttt{DatasetMetric}.
    \item Generic \texttt{MI} must implement trait \texttt{UnboundedMetric}.
    \item Generic \texttt{MO} must implement trait \texttt{ApproximateMeasure}.
\end{itemize}

\subsubsection*{User-verified}
None

\subsection*{Pseudocode}
\lstinputlisting[language=Python,firstline=2,escapechar=|]{./pseudocode/make_private_group_by.py}

\subsubsection*{Postconditions}
\begin{theorem}
    \label{postcondition}
    \validMeasurement{\texttt{(input\_domain, input\_metric, \\output\_measure, plan, global\_scale, threshold, MS, MI, MO)}}{\texttt{make\_private\_group\_by}}
\end{theorem}
\section{Proof}

We now prove the postcondition (Theorem~\ref{postcondition}).
\begin{proof}

The function logic breaks down into parts:
\begin{enumerate}
    \item establish stability of input (line \ref{line:input_stability})
    \item prepare for release of \texttt{aggs} (line \ref{line:prep-release-aggs})
    \item prepare for release of \texttt{keys} (line \ref{line:prep-release-keys})
    \begin{enumerate}
        \item reconcile information about the threshold (line \ref{line:reconcile-threshold})
        \item update key sanitizer (line \ref{line:final-sanitizer})
    \end{enumerate}
    \item build final measurement (line \ref{line:build-meas})
    \begin{enumerate}
        \item construct function (line \ref{line:function})
        \item construct privacy map (line \ref{line:privacy-map})
    \end{enumerate}
\end{enumerate}

\rustdoc{measurements/make\_private\_lazyframe/group\_by/matching/fn}{match\_group\_by} returns a valid
\rustdoc{measurements/make\_private\_lazyframe/group\_by/matching/struct}{MatchGroupBy} struct. 
In this struct, \texttt{input} is the input plan, \texttt{keys} is the grouping keys, \texttt{aggs} is the list of expressions to compute per-partition,
and \texttt{key\_sanitizer} details how to sanitize the key-set (line \ref{line:match-group-by}).

\subsection{Stability of input}
Start by establishing properties of the following variables, 
which hold for any setting of the input arguments.

By the postcondition of \texttt{StableDslPlan.make\_stable}, \texttt{t\_prior} is a valid transformation (line \ref{line:tprior}).
By the postcondition of \texttt{match\_grouping\_columns}, \texttt{grouping\_columns} holds the names of the grouping columns.
\texttt{margin} denotes what is considered public information about the key set, 
pulled from descriptors in the input domain (line \ref{line:groupcols}).

When sanitizing keys via a join, an upper bound on the total number of partitions
can be statically derived via the length of the grouping keys.
Line \ref{line:keys-len} retrieves this information from the key-set and assigns it to the margin.

\texttt{is\_join} indicates that key sanitization will occur via a join. 
This will be a useful criteria later.

\subsection{Prepare to release \texttt{aggs}}
Line \ref{line:joint} starts the process to prepare a joint measurement 
for releasing the per-partition aggregations \texttt{aggs}.
Each measurement's input domain is the wildcard expression domain,
used to prepare computations that will be applied over data grouped by grouping columns \texttt{by}.

By the postcondition of \texttt{make\_basic\_composition}, 
\texttt{m\_exprs} is a valid measurement that prepares a batch of expressions that, 
when executed via \texttt{f\_comp}, satisfies the privacy guarantee of \texttt{f\_privacy\_map}.

Now that we've prepared the necessary prerequisites for privatizing the aggregations,
we switch to privatizing the keys.

\subsection{Prepare to release \texttt{keys}}
\texttt{key\_sanitization} needs to be updated with information that was not available in the initial match on line \ref{line:match-group-by}.
\begin{itemize}
    \item When filtering, a threshold may be passed into the constructor, and we must determine a suitable column to filter/threshold against.
    \item When joining, we need expressions for filling null values corresponding to partitions that don't exist in the sensitive data.
\end{itemize}


We first reconcile information about the filtering threshold (line \ref{line:reconcile-threshold}).
\begin{itemize}
    \item In the setting where grouping keys are considered public, 
    or key sanitization is handled via a join, no thresholding is necessary (line \ref{line:no-thresholding}).
    \item Otherwise, if the key sanitizer contains filtering criteria (line \ref{line:sanitizer-threshold}),
    then by the postcondition of \texttt{find\_len\_expr}, filtering on \texttt{name} can be used to satisfy 
    $\delta$-approximate DP.
    \texttt{noise} of type \rustdoc{measurements/make\_private\_expr/fn}{NoisePlugin} details the noise distribution and scale. 
    \texttt{threshold\_info} then contains the column name, noise distribution, threshold value and whether a filter needs to be inserted into the query plan.
    In this case, since the threshold comes from the query plan, it is not necessary to add it to the query plan, and is therefore false.
    \item In the case that a threshold has been provided to the constructor (line \ref{line:constructor-threshold}),
    then \texttt{find\_len\_expr} will search for a suitable column to threshold on, 
    returning with the \texttt{name} and \texttt{noise} distribution of the column.
    Since the threshold comes from the constructor and not the plan, it will be necessary to add this filtering threshold to the query plan (explaining the true value).
    \item By line \ref{line:not-private} no suitable filtering criteria have been found, 
    and by the first case there is no suitable invariant for the margin or explicit join keys,
    so it is not possible to release the keys in a way that satisfies differential privacy,
    and the constructor refuses to build a measurement.
\end{itemize}

In common use through the context API, if a mechanism is allotted a delta parameter for stable key release but doesn't already satisfy approximate-DP,
then a search is conducted for the smallest suitable threshold parameter.
The branching logic from line \ref{line:reconcile-threshold} is intentionally written to ignore the constructor threshold 
when a suitable filtering threshold is already detected in the plan, to avoid overwriting/changing it.

We now update \texttt{key\_sanitization} starting from line \ref{line:final-sanitizer}:
\begin{itemize}
    \item When filtering (line \ref{line:incorporate-threshold}), \texttt{threshold\_info} will always be set. 
    \texttt{threshold\_expr} reflects the reconciled criteria, using the chosen filtering column and threshold. 
    This threshold expression is applied either way the logic branches on line \ref{line:new-threshold-sanitizer}.
    The first case preserves any additional filtering criteria that was already present in the plan, but not used for key release.
    \item When joining (line \ref{line:incorporate-join-fill}) the sanitizer needs a way to fill missing values from partitions missing in the data.
    This is provided by \texttt{null\_exprs}, which contain imputation strategies for filling in missing values 
    in a way that is indistinguishable from running the mechanism on an empty partition.
\end{itemize}

\texttt{key\_sanitizer} now contains all necessary information to ensure that the keys are sanitized, 
and will be used to construct the function.
\texttt{threshold\_info}, \texttt{is\_join} and \texttt{public\_info} are consistent with \texttt{key\_sanitizer}, 
and will be used to construct the privacy map.

\subsection{Build final measurement}
We now move on to the implementation of the function on line \ref{line:function}, 
using all of the properties shown of the variables established thus far.
The function returns a DslPlan that applies each expression from \texttt{m\_exprs}
to \texttt{arg} grouped by \texttt{keys}.
\texttt{key\_sanitizer} is conveyed into the plan, if set, 
to ensure that the keys are also privatized if necessary.

In the case of the join privatization, by the definition of \texttt{KeySanitizer},
the join will either be a left or right join.
The branching swaps the input plan and labels plan to ensure that the sensitive input data is always joined against the labels,
but using the same join type as in the original plan.
Once the join is applied, the fill imputation expressions are applied, hiding which partitions don't exist in the original data.

It is assumed that the emitted DSL is executed in the same fashion as is done by Polars.
This proof/implementation does not take into consideration side-channels involved in the execution of the DSL.

We now move on to the implementation of the privacy map. \ref{line:privacy-map}
The measurement for each expression expects data set distances in terms of a triple:
\begin{itemize}
    \item $L^0$: the greatest number of partitions that can be influenced by any one individual. 
    This is no greater than the input distance (an individual can only ever influence as many partitions as they contribute rows), 
    but could be smaller when supplemented by 
    the \texttt{max\_influenced\_partitions} metric descriptor or \texttt{max\_num\_partitions} domain descriptor.
    \item $L^\infty$: the greatest number of records that can be added or removed by any one individual in each partition. 
    This is no greater than the input distance, 
    but could be tighter when supplemented by 
    the \texttt{max\_partition\_contributions} metric descriptor or the \texttt{max\_partition\_length} domain descriptor.
    \item $L^1$: the greatest total number of records that can be added or removed across all partitions.
    This is no greater than the input distance,
    but could be tighter when accounting for the $L^0$ and $L^\infty$ distances.
\end{itemize}

By the postcondition of \texttt{f\_privacy\_map}, the privacy loss of releasing the output of \texttt{aggs}, 
when grouped data sets may differ by this distance triple,
is \texttt{d\_out}.

We also need to consider the privacy loss from releasing \texttt{keys}.
On line \ref{line:privacy-map-static-keys} under the \texttt{public\_info} invariant, or under the join sanitization, 
releases on any neighboring datasets $x$ and $x'$ will share the same key-set,
resulting in zero privacy loss.

% If the grouping keys are private, then the threshold has been prepared in \ref{reconcile-threshold}.
% At this point, it is necessary to show that the probability that any unstable partition is released is at most $\delta$.
% An unstable partition is any partition that may not exist on a neighboring data set:
% that is, an unstable partition is a partition whose data is unique to an individual.
% Therefore, unstable partitions have \texttt{li} records or fewer.

% The distance to instability (\texttt{d\_instability}) denotes the gap between the size of the largest unstable partition and the threshold.
% The additional delta parameter (\texttt{delta\_p}) denotes the worst-case probability 
% that an unstable partition an individual contributes is still present in the release,
% because noise added to release the DP count estimate exceeds \texttt{d\_instability}.

We now adapt the proof from \cite{rogers2023unifyingprivacyanalysisframework} (Theorem 7) 
to consider the case of stable key release from line \ref{line:privacy-map-threshold}.
Consider $S$ to be the set of labels that are common between $x$ and $x'$.
Define event $E$ to be any potential outcome of the mechanism for which all labels are in $S$
(where only stable partitions are released).
We then lower bound the probability of the mechanism returning an event $E$.
In the following, $c_j$ denotes the exact count for partition $j$,
and $Z_j$ is a random variable distributed according to the distribution used to release a noisy count.

\begin{align*}
    \Pr[E] &= \prod_{j \in x \backslash x'} \Pr[c_j + Z_j \le T] \\
    &\ge \prod_{j \in x \backslash x'} \Pr[\Delta_\infty + Z_j \le T] \\
    &\ge \Pr[\Delta_\infty + Z_j \le T]^{\Delta_0}
\end{align*}

The probability of returning a set of stable partitions ($\Pr[E]$) 
is the probability of not returning any of the unstable partitions.
We now solve for the choice of threshold $T$ such that $\Pr[E] \ge 1 - \delta$.

\begin{align*}
    \Pr[\Delta_\infty + Z_j \le T]^{\Delta_0} &= \Pr[Z_j \le T - \Delta_\infty]^{\Delta_0} \\
    &= (1 - \Pr[Z_j > T - \Delta_\infty])^{\Delta_0}
\end{align*}

Let \texttt{d\_instability} denote the distance to instability of $T - \Delta_\infty$.
By the postcondition of \\ \texttt{integrate\_discrete\_noise\_tail},
the probability that a random noise sample exceeds \texttt{d\_instability} is at most \texttt{delta\_single}.
Therefore $\delta = 1 - (1 - \texttt{delta\_single})^{\Delta_0}$.
This privacy loss is then added to \texttt{d\_out}.

Together with the potential increase in delta for the release of the key set,
then it is shown that \function(x), \function(x') are \dout-close under \texttt{output\_measure}.

\end{proof}

\bibliographystyle{alpha}
\bibliography{ref}
\end{document}